\section{Introduction}\label{s:intro}
University rocketry teams such as RISE~\cite{rise} compete in events like the Spaceport America Cup and the European Rocketry Challenge (EuRoC). There they launch high-powered rockets (HPR) to altitudes of up to 9~km, carrying scientific payloads~\cite{rocketpy-ref}. The rockets run flight software (FSW) that processes sensor data, detects flight phases, deploys parachutes and, on more advanced vehicles, controls air brakes~\cite{rocketpy-ref, clary2024georgia-tech-avionics}. Apogee detection, the moment the rocket reaches its peak altitude, is the critical case: an early or missed parachute deployment cannot be corrected in flight. As missions grow more ambitious, the software grows more complex, and so does the risk that a fault goes unnoticed until launch day.

Testing against live launches is impractical. Launches are expensive, uncontrolled and irreversible~\cite{coopmans-sil-hil-verification-uav, kiesbye2019-moveII-cubsat}, and annual competitions offer no second launch window, so a failed flight can ruin a team's season. UAV and small-spacecraft programmes instead close a feedback loop between the FSW and a simulated model of the vehicle and its environment, the \emph{plant}. The FSW runs either on the host (software-in-the-loop, SIL) or on the target hardware (hardware-in-the-loop, HIL)~\cite{santos-hil-uav-autopilot, bittar-guidance-sil-uav-xplane-simulink, kiesbye2019-moveII-cubsat}. Institutional programmes such as ESA and NASA have built such infrastructure over years of investment. Student teams, with volunteer labour, small budgets and single-use rockets, have not. They rely on open-source tools instead: OpenRocket~\cite{niskanen2013openrocket} and RocketPy~\cite{rocketpy-ref} for trajectories, and NASA's F Prime~\cite{bocchino2018fprime} for flight software. No framework couples them.

Existing rocket simulators are open-loop. OpenRocket, the Cambridge Rocketry Simulator and RocketPy~\cite{niskanen2013openrocket, eerland2017cambridgerocketry, rocketpy-ref} predict a trajectory offline from geometry, thrust curve and environment. None feeds live sensor readings to running software or accepts its control outputs, so teams are left with unit tests and open-loop replay, which cannot show how the FSW's decisions change the flight. None can inject sensor faults either, so failure-recovery logic stays untested until flight.

This paper presents RITL (Rocket in the Loop), which closes that loop between RocketPy and F Prime through a Python Orchestrator. RocketPy is a continuous-time simulator and F Prime is discrete and event-driven, so coupling them is a heterogeneous co-simulation problem~\cite{gomes2017cosimulation}. The Orchestrator offers three coupling strategies that trade synchronisation against timing realism and wall-clock cost, runs the same FSW in SIL and HIL, and injects faults inside the loop. We address three research questions (RQ):
\begin{itemize}
\item \textbf{RQ1:} How can a closed-loop environment for HPR flight software be built from RocketPy and F Prime? We give an architecture that exchanges sensor and actuation data at every controller step without modifying RocketPy's solver, and runs the same FSW in SIL and HIL with a recompilation only.
\item \textbf{RQ2:} How do the coupling strategy and the execution context (SIL or HIL) affect apogee accuracy, parachute timing and wall-clock cost? We characterise three strategies, Lockstep, Snapshot and Rate Group, in SIL and on a Raspberry Pi 4.
\item \textbf{RQ3:} Can the framework inject a sensor fault during flight, and does the FSW recover so that parachute deployment still succeeds? We demonstrate in-loop fault injection with a barometer freeze.
\end{itemize}

RQ1 is answered by a working implementation, not a measurable hypothesis. RQ2 needs an accuracy criterion, so we take an apogee error below 12\,m as the acceptance threshold for pre-flight verification. It comes from the Camões rocket configuration used throughout~\cite{rocketpy_camoes}. Camões, built by Rocket Experiment Division, reached 3015\,m at EuRoC 2023, while the standalone RocketPy prediction was 3003\,m, 12\,m lower. A coupling error below that existing prediction uncertainty does not reduce the value of the simulation compared with using RocketPy alone. RQ3 matters because a sensor fault can prevent correct apogee detection and parachute deployment, an irreversible failure on a live launch. Testing the recovery logic in simulation first reduces that risk.

In short, Lockstep reproduces a standalone RocketPy baseline within 0.08\,m, whereas the non-blocking Snapshot and Rate Group strategies add a one-tick actuation delay and overshoot apogee by about 4\,m and 3.84--7.39\,m on average. All mean errors stay within the 12\,m threshold. After an injected barometer freeze, the FSW deployed the drogue parachute within 0.1\,s of true apogee.

Section~\ref{s:background} gives background and Section~\ref{s:related} related work. Section~\ref{closed-loop-sim-env} presents the framework (RQ1), Section~\ref{experiment-design} the experiments (RQ2, RQ3) and Section~\ref{s:evaluation} their results. Section~\ref{discussion} discusses implications and threats to validity, and Section~\ref{s:conclusion} concludes.
